\chapter{Kontsevich--Segal Functorial Quantum Field Theory}
\label{ch:kontsevich-segal-functorial-qft}
\ConventionsEuclidean

\section{Functorial Data and the Path-Integral Problem}

The Wightman, Osterwalder--Schrader, and Haag--Kastler frameworks encode
different shadows of locality.  Wightman fields encode distributional
Lorentzian local operators on a Hilbert space; OS data encode Euclidean
reflection-positive correlation distributions; Haag--Kastler nets encode
bounded observable algebras localized in spacetime regions.  The
Kontsevich--Segal viewpoint begins from the geometric operation that a path
integral is expected to satisfy before any choice of coordinates on fields is
made: cut a background along hypersurfaces, assign state spaces to the cuts,
and recover the amplitude of the glued background by composing the amplitudes
of the pieces.

Let \(D\) be the spacetime dimension.  The intended datum is a symmetric
monoidal assignment
\[
  Z:\operatorname{Bord}_{D}^{\mathrm{adm}}\longrightarrow
  \operatorname{TVect}_{\C},
\]
where \(\operatorname{Bord}_{D}^{\mathrm{adm}}\) is a category of
\((D-1)\)-manifolds and \(D\)-dimensional bordisms with admissible complex
metrics and auxiliary background fields, and
\(\operatorname{TVect}_{\C}\) is a category of complex topological vector
spaces and continuous linear maps.  The target cannot be restricted in
advance to Hilbert spaces.  Hilbert spaces arise from a reflection-positive
pairing on suitable real slices; before that pairing is imposed and
completed, the natural state spaces are locally convex spaces or distribution
spaces.

The word ``admissible'' carries the physical content of positive energy.  The
complex metric is allowed to move inside a domain containing Riemannian
metrics and having Lorentzian metrics on its boundary.  Analytic dependence
inside this domain replaces the upper-half-plane analyticity of
\(\exp(\ii tH)\) in quantum mechanics.  The objective is therefore sharper
than a generic functorial field theory: the functor must be holomorphic in
admissible complex metrics, compatible with sewing, and equipped with a
reflection structure whose Lorentzian boundary values produce unitary
evolution and microcausal fields when the required analytic hypotheses hold.

\section{Segal Sewing in Two-Dimensional CFT}

The two-dimensional conformal case is the prototype.  Let
\(\operatorname{Bord}^{\mathrm{conf}}_2\) denote the category whose objects
are finite disjoint unions of analytically parametrized oriented circles and
whose morphisms are compact Riemann surfaces with incoming and outgoing
parametrized boundary components.  Composition is sewing along boundary
circles with inverse parametrizations.  A Segal CFT assigns a topological
vector space \(\Hilb_{S}\) to each object \(S\) and a continuous linear map
\[
  Z(\Sigma):\Hilb_{S_{\mathrm{in}}}\longrightarrow
  \Hilb_{S_{\mathrm{out}}}
\]
to a Riemann surface \(\Sigma:S_{\mathrm{in}}\to S_{\mathrm{out}}\), with
\[
  Z(\Sigma_2\circ\Sigma_1)=Z(\Sigma_2)Z(\Sigma_1),
  \qquad
  Z(\Sigma\sqcup\Sigma')=Z(\Sigma)\otimes Z(\Sigma').
\]
The sewing map is required to be continuous or holomorphic with respect to
the relevant moduli of surfaces and boundary parametrizations.  Reflection of
a Riemann surface across a parametrized boundary circle gives an antilinear
involution on the state space, and reflection positivity requires the doubled
surface to define a positive semidefinite pairing.

The operator product in a Segal CFT is a sewing statement.  A pair of
punctures contained in a disk determines a state on the boundary circle of
that disk.  Expanding this state in the spectral decomposition of the annulus
semigroup gives the local fields that may replace the pair of insertions
when the disk is sewn into a larger surface.  Thus convergence of the OPE is
not an extra mnemonic; it is the convergence of a state-vector expansion in
the boundary state space chosen by the sewing theory.

\section{Allowable Complex Metrics}

Let \(V\) be a real vector space of dimension \(D\).  A complex metric on
\(V\) is a nondegenerate complex-valued symmetric bilinear form
\[
  g\in\operatorname{Sym}^2(V^*)\otimes_{\R}\C .
\]
Let \(|\wedge^D V^*|\) be the real density line.  A choice of square root of
\(\det g\) defines a complex density
\[
  \operatorname{vol}_g=(\det g)^{1/2}|dy^1\cdots dy^D|
  \in |\wedge^D V^*|\otimes_{\R}\C
\]
in a basis \(y^1,\ldots,y^D\).  The sign of the square root is fixed by
requiring the \(0\)-form positivity condition below.

\begin{definition}[Kontsevich--Segal allowability]
\label{def:ks-allowable-linear}
A complex metric \(g\) on \(V\) is allowable when, for every
\(q=0,\ldots,D\), the quadratic map
\[
  Q_{q,g}:\wedge^q V^*\longrightarrow |\wedge^D V^*|\otimes_{\R}\C,
  \qquad
  \alpha\longmapsto \alpha\wedge *_g\alpha
\]
has positive real part on every nonzero real \(q\)-form \(\alpha\).  For
\(q=0\), this says that \(\operatorname{Re}\operatorname{vol}_g\) is a
positive real density.
\end{definition}

For a \(q\)-form gauge field strength \(F\), the Euclidean quadratic action is
\[
  I_q(F;g)
  =
  \frac{1}{2q!}
  \int_M
  \sqrt{\det g}\,
  g^{i_1j_1}\cdots g^{i_qj_q}
  F_{i_1\cdots i_q}F_{j_1\cdots j_q}\,\dd^Dx .
\]
Definition~\ref{def:ks-allowable-linear} is precisely the pointwise
condition that the real part of this quadratic form be positive for every
real \(q\)-form field strength, simultaneously for all \(q\).  Imposing all
degrees is essential because a field theory may contain scalars, gauge
fields, higher-form fields, dual fields, and quantized flux sectors whose
integration cycles cannot all be repaired by a degree-dependent phase
rotation.

\begin{theorem}[Argument criterion for allowability]
\label{thm:ks-allowability-argument}
Let \(g\) be a complex metric on \(V\).  It is allowable if and only if there
is a real basis in which
\[
  g=\sum_{a=1}^D \lambda_a (dy^a)^2,
  \qquad
  \lambda_a\in\C\setminus(-\infty,0],
\]
and, with the principal argument \(\theta_a=\operatorname{Arg}\lambda_a\in
(-\pi,\pi)\),
\begin{equation}
  \sum_{a=1}^D |\theta_a|<\pi .
  \label{eq:ks-allowability-argument}
\end{equation}
\end{theorem}

\begin{proof}
Assume first that \(g\) is allowable.  The \(q=0\) condition chooses the
branch of \((\det g)^{1/2}\) with positive real part.  The \(q=1\) condition
states that the complex symmetric matrix
\[
  W^{ij}=(\det g)^{1/2}g^{ij}
\]
has positive-definite real part.  Write \(W=A+\ii B\) with \(A,B\) real
symmetric and \(A>0\).  A real change of basis puts \(A\) in the identity
form.  A further orthogonal change of basis diagonalizes \(B\).  Hence \(W\)
is diagonal in a real basis, so \(W^{-1}\) and then \(g\) are diagonal in the
same basis.  Write the diagonal entries as \(\lambda_a\).

For a subset \(S\subset\{1,\ldots,D\}\), the real basis \(q\)-form
\(\dd y^S=\wedge_{a\in S}\dd y^a\) has
\begin{equation}
  \dd y^S\wedge *_g\dd y^S
  =
  A_S\,|dy^1\cdots dy^D|,
  \qquad
  A_S=(\det g)^{1/2}\prod_{a\in S}\lambda_a^{-1}.
  \label{eq:ks-form-coefficient}
\end{equation}
The \(q\)-form positivity conditions are therefore equivalent to
\(\operatorname{Re}A_S>0\) for all subsets \(S\).  If
\(\theta_a=\operatorname{Arg}\lambda_a\), then
\[
  \operatorname{Arg} A_S
  =
  \frac12\sum_{a=1}^D\theta_a-\sum_{a\in S}\theta_a
  =
  \frac12\left(\sum_{a\notin S}\theta_a-\sum_{a\in S}\theta_a\right).
\]
Let \(P=\{a:\theta_a\ge0\}\).  Positivity for \(S=P\) gives
\[
  -\frac{\pi}{2}
  <
  -\frac12\sum_{a=1}^D|\theta_a|
  <
  \frac{\pi}{2},
\]
and hence \eqref{eq:ks-allowability-argument}.  Conversely, if
\eqref{eq:ks-allowability-argument} holds, then for every subset \(S\),
\[
  |\operatorname{Arg}A_S|
  \le
  \frac12\sum_{a=1}^D|\theta_a|
  <
  \frac{\pi}{2}.
\]
Thus every \(A_S\) has positive real part.  Since the form coefficients are
diagonal in the real basis, every nonzero real \(q\)-form has
\(\operatorname{Re}Q_{q,g}>0\).  This is
Definition~\ref{def:ks-allowable-linear}.
\end{proof}

\begin{example}[Euclidean, Lorentzian, and multi-time signatures]
\label{ex:ks-basic-metrics}
A positive Riemannian metric has all \(\theta_a=0\), hence is allowable.  A
Lorentzian metric with diagonal entries \((-1,1,\ldots,1)\) has one argument
equal to \(\pi\), so it lies on the boundary of the allowable domain.  The
regularized metrics
\[
  -(1-\ii\epsilon)(dy^1)^2+\sum_{a=2}^D(dy^a)^2,
  \qquad
  -(1+\ii\epsilon)(dy^1)^2+\sum_{a=2}^D(dy^a)^2,
  \qquad \epsilon>0,
\]
lie in the two adjacent allowable components approaching that Lorentzian
boundary point.  A real metric with two negative directions has two arguments
equal to \(\pi\) and does not lie on this boundary.  The criterion therefore
distinguishes Lorentzian signature from signatures with more than one time
direction.
\end{example}

\section{The Admissible Bordism Category}

An admissible \(D\)-bordism is a smooth compact \(D\)-manifold \(M\) with
boundary decomposed into incoming and outgoing collars, together with:
\[
  g\in\Gamma(\operatorname{Sym}^2T^*M\otimes\C),
  \qquad
  \text{background bundles and sources},
\]
where \(g_x\) is allowable for every \(x\in M\).  Boundary collars are part of
the datum so that gluing is defined by identifying collars with matching
induced complex metrics and background fields.  The category is symmetric
monoidal under disjoint union.

The state space assigned to a boundary germ \(Y\) will be denoted
\[
  E(Y).
\]
The notation emphasizes that \(Y\) is a germ of a hypersurface together with
its collar metric and background fields, rather than only the underlying
\((D-1)\)-manifold.  For a bordism \(M:Y_0\to Y_1\), the amplitude is a
continuous linear map
\[
  Z(M):E(Y_0)\longrightarrow E(Y_1).
\]
For a closed \(M\), \(Z(M)\in\C\) is the partition function.  For marked
points or extended insertions, \(E(Y)\) is replaced by a multilinear
observable bundle over the configuration space of insertions.

\begin{definition}[K-S functorial QFT datum]
\label{def:ks-functorial-qft}
A Kontsevich--Segal functorial QFT on a specified class of backgrounds
consists of:
\begin{enumerate}
\item a symmetric monoidal functor \(Z\) from admissible complex bordism
categories to complex topological vector spaces;
\item holomorphic dependence of \(Z(M,g)\) on the admissible complex metric
and on complexified sources in local coordinate charts on the space of
background data;
\item collar locality: if \(M=M_2\circ_Y M_1\), then
\[
  Z(M)=Z(M_2)Z(M_1)
\]
as a continuous map after the matching state-space identifications along
\(Y\);
\item a reflection operation \(Y\mapsto \bar Y\) and \(M\mapsto \bar M\)
compatible with complex conjugation of metric and sources;
\item reflection positivity on real Riemannian reflection doubles, as stated
below.
\end{enumerate}
\end{definition}

\section{Reflection Positivity and Hilbert Spaces}

Let \(Y\) be a real Riemannian boundary germ fixed by an orientation-reversing
reflection \(\rho\).  Suppose \(M_+\) is a bordism from the empty object to
\(Y\), and let \(\bar M_+\) be its reflected conjugate from \(Y\) to the empty
object.  The double
\[
  \bar M_+\circ_Y M_+
\]
is a closed Riemannian background.  If \(u,v\in E(Y)\) arise from states
created by such half-bordisms with insertions, define
\[
  \langle u,v\rangle_Y
  =
  Z(\bar M_u\circ_Y M_v).
\]
Reflection positivity requires this sesquilinear form to be positive
semidefinite on the subspace generated by positive-side states.  The physical
Hilbert space associated to \(Y\) is then
\begin{equation}
  \Hilb_Y
  =
  \overline{E_+(Y)/\{u:\langle u,u\rangle_Y=0\}},
  \label{eq:ks-hilbert-completion}
\end{equation}
where the bar denotes Hilbert completion.

In a real-analytic globally hyperbolic Lorentzian spacetime, a Lorentzian
time-slab appears as a boundary limit of admissible complex slabs.  Under the
holomorphic and boundedness hypotheses needed to take this limit, the K-S
construction gives unitary propagation between the Hilbert spaces of
time-symmetric hypersurfaces.  Local field operators on a time slice are
obtained as limits of insertion bordisms pinched to the slice; operators at
spacelike separated points commute because both insertions can be represented
on one common slice and the Euclidean sewing map is symmetric under exchanging
their order.

\section{Comparison with OS, Wightman, and Local Nets}

\begin{proposition}[K-S data give OS Schwinger functions on flat space]
\label{prop:ks-to-os}
Assume a K-S functorial QFT is defined on Euclidean \(\R^D\) with
translation-invariant vacuum state, local observable bundles
\(\mathcal O_x\), and tempered boundary values for point insertions.  For
local observables \(A_1,\ldots,A_n\), define
\[
  S_n(A_1,x_1;\ldots;A_n,x_n)
  =
  Z(\R^D;A_1(x_1),\ldots,A_n(x_n))
\]
for distinct Euclidean points, extended as a distribution by the assumed
tempered boundary value.  Then Euclidean covariance, symmetry at separated
points, reflection positivity, and clustering of the Schwinger functions
follow from the functorial axioms plus the corresponding covariance and
vacuum uniqueness assumptions on \(Z\).
\end{proposition}

\begin{proof}
Euclidean covariance is functoriality under background isomorphisms.
Symmetry follows by moving separated point insertions inside a Riemannian
background without crossing; the configuration-space complement of diagonals
records precisely the possible permutation monodromy, with ordinary bosonic
local observables having trivial monodromy.  Reflection positivity is the
positivity of the doubled half-space bordism.  Clustering is not a pure sewing
identity: it is the assertion that the translation semigroup on a long
cylinder has a unique zero-energy vector and that the remaining spectrum is
separated or decays in the topology used for the state spaces.  Under that
spectral hypothesis, inserting a long cylinder between two groups of
observables projects to the vacuum state, giving the OS cluster property.
\end{proof}

\begin{quotedtheorem}[K-S Lorentzian boundary construction]
\label{thm:ks-lorentzian-boundary}
For a K-S theory on admissible complex metrics with the analytic dependence,
reflection structure, and boundedness assumptions required by
Kontsevich--Segal, real-analytic globally hyperbolic Lorentzian bordisms are
obtained as boundary values of admissible complex bordisms.  Time-symmetric
hypersurfaces acquire Hilbert spaces, globally hyperbolic Lorentzian
bordisms induce unitary maps between them, and local fields obtained from
insertion germs satisfy spacelike commutativity.  This is the content of the
Lorentzian-boundary construction in Kontsevich--Segal's
\emph{Wick Rotation and the Positivity of Energy in Quantum Field Theory},
with Witten's \emph{A Note on Complex Spacetime Metrics} giving a detailed
analysis of the finite-dimensional allowability criterion and gravitational
examples.
\end{quotedtheorem}

Combining Proposition~\ref{prop:ks-to-os} with OS reconstruction gives the
route
\[
  \text{K-S functorial data}
  \longrightarrow
  \text{OS Schwinger functions}
  \longrightarrow
  \text{Wightman fields}.
\]
After Wightman fields have been reconstructed, local nets are obtained by
taking von Neumann algebras generated by bounded functions of smeared fields,
provided the self-adjointness and strong-locality hypotheses needed for those
bounded functions are verified.  Thus the arrow from K-S to Haag--Kastler is
a composite arrow with analytic and operator-domain assumptions at each
stage.

The reverse direction requires additional structure.  OS Schwinger functions
on flat space do not by themselves assign state spaces to arbitrary
hypersurface germs, nor do they define amplitudes for all admissible complex
bordisms.  A reconstruction of a K-S functor from OS data therefore requires
sewing estimates, collar factorization, and control of metric dependence.
Those hypotheses are available in selected two-dimensional constructions and
topological theories; they are open in physically central higher-dimensional
interacting theories.

\section{OPE from Sewing}

Let \(B_r(x)\subset M\) be a small ball whose boundary sphere is \(S_r(x)\).
A collection of local insertions inside \(B_r(x)\) defines, by the K-S
amplitude of the punctured ball, a vector
\[
  \Psi_{A_1,\ldots,A_k}(r;\xi_1,\ldots,\xi_k)\in E(S_r(x)),
\]
where \(\xi_i\) are coordinates of the insertion points relative to \(x\).
Suppose the annular evolution identifying \(E(S_r(x))\) with \(E(S_1(x))\)
has a discrete generalized eigenspace decomposition
\[
  E(S_1(x))=\widehat{\bigoplus_{\Delta,s}} E_{\Delta,s}
\]
with nuclear convergence for the expansion of punctured-ball states.  Choose
a basis of local operators \(\mathcal O_{\Delta,s,a}\) corresponding to the
dual basis of these eigenspaces.  Then sewing the punctured ball into the
complement \(M\setminus B_r(x)\) gives
\begin{equation}
  A_1(x+\xi_1)\cdots A_k(x+\xi_k)
  \sim
  \sum_{\Delta,s,a}
  C^{\Delta,s,a}_{A_1\ldots A_k}(\xi_1,\ldots,\xi_k)\,
  \mathcal O_{\Delta,s,a}(x),
  \label{eq:ks-ope-from-sewing}
\end{equation}
with convergence in the topology in which the boundary-state expansion
converges.  In a CFT, annular evolution is generated by dilations, so the
coefficients have homogeneous leading behavior determined by scaling
dimensions and spin.  In a massive theory the same construction gives a local
short-distance expansion when the small-ball state admits the corresponding
asymptotic spectral expansion.

This derivation identifies the analytic hypotheses behind OPE convergence:
state spaces on spheres, trace or nuclear bounds for annuli, spectral
control of the annular semigroup, and continuity of sewing.  Without these
hypotheses, \eqref{eq:ks-ope-from-sewing} is a formal expansion rather than a
theorem.

\section{Stoyanovsky-Style Wick Rotation}

Stoyanovsky's axiomatic Wick-rotation program may be viewed as a
hypersurface-evolution formulation of the same problem.  The basic object is
a bundle of state spaces over a space of embedded hypersurfaces in a
complexified spacetime, equipped with flat evolution maps for deformations of
the hypersurface.  Local fields act by insertion when the hypersurface passes
the insertion point.  Flatness expresses independence of the chosen
interpolating family of hypersurfaces; positivity and boundary-value
conditions select the Lorentzian real slice.

The K-S metric formulation supplies a concrete admissible domain for those
hypersurface deformations.  Rather than complexifying the time coordinate
alone, it complexifies the metric and imposes
\eqref{eq:ks-allowability-argument} pointwise.  This gives a finite-dimensional
linear-algebra condition at each tangent space and a gluing category of
admissible metric bordisms.  The two viewpoints should agree when both are
defined: hypersurface evolution is the local expression of the functor on
thin bordisms.

\section{Examples and Construction Status}

Topological theories are the cleanest functorial examples.  If \(Z\) is
independent of the metric and satisfies Atiyah--Segal sewing, then it is a
K-S-type theory restricted to the topological subcategory: the admissible
metric dependence is constant.  Finite gauge theory, Dijkgraaf--Witten
theory, and Reshetikhin--Turaev theories provide theorem-level examples in
this topological sector, with their anomaly and extended-boundary data
recorded in Volume~VIII.

Two-dimensional CFT supplies the deepest non-topological testing ground.
Segal's sewing axioms have been verified in important representation-theoretic
and probabilistic settings, including free-fermionic and loop-group examples
in the operator-algebraic literature and Liouville-theory sewing results in
the probabilistic conformal-bootstrap program.  Each such result is a
two-dimensional Segal-CFT theorem with its own topology on state spaces and
moduli dependence.  Passing from these theorems to a general-dimensional K-S
construction is a separate problem.

Two-dimensional Yang--Mills has exact heat-kernel formulas on surfaces and a
rigorous measure-theoretic construction in several formulations.  These
formulas obey cutting identities at the level of class functions and
holonomy kernels.  A full continuum Segal functor with all state spaces,
collar topologies, and boundary field insertions specified is a further
structure that must be supplied before calling it a completed K-S example.

Free massive scalar theory in \(D\ge3\) illustrates the gap.  The Gaussian
path integral gives explicit kernels on cylinders, and the allowability
condition gives positivity of the real part of the quadratic action.  A
complete K-S construction still has to specify state spaces for arbitrary
hypersurface germs, prove trace or nuclear sewing for admissible collars,
control point-field insertions as distributions or hyperfunctions, prove
reflection positivity for the chosen completions, and show holomorphic
dependence on the admissible metric.  The monograph therefore treats the free
scalar in \(D\ge3\) as a principal construction problem rather than as an
already completed K-S theorem.

\section{Research Ledger}

\begin{openproblem}[Free fields as K-S functors]
Construct the full K-S functor for the free massive scalar field in
\(D\ge3\).  The construction must specify hypersurface state spaces,
Gaussian amplitudes for admissible collars, holomorphic metric dependence,
reflection positivity, distributional field insertions, and sewing estimates.
\end{openproblem}

\begin{openproblem}[\(P(\phi)_2\) and constructive scalar models]
Starting from constructive OS data for \(P(\phi)_2\), build state spaces on
curves, prove collar factorization, and extend the Schwinger functions to a
Segal/K-S functor on the relevant admissible two-dimensional bordism
category.
\end{openproblem}

\begin{openproblem}[Liouville and nonrational CFT]
Determine the exact relation between probabilistic Liouville correlators,
the DOZZ/bootstrap construction, Segal sewing, and a complete functorial
state-space assignment.  The required output is a functor with specified
boundary state spaces and sewing maps, not only closed-surface correlators.
\end{openproblem}

\begin{openproblem}[Two-dimensional Yang--Mills]
Upgrade the heat-kernel and holonomy-cutting formulas of two-dimensional
Yang--Mills to a continuum Segal functor with boundary Hilbert spaces,
defect insertions, and compatibility with the admissible complex metric
domain.
\end{openproblem}

\begin{openproblem}[Three-dimensional constructive models]
For \(\phi^4_3\), Gross--Neveu type models, and related constructive
theories, identify the additional estimates needed to pass from OS
correlators to state spaces on arbitrary surfaces and to admissible
complex-metric sewing.
\end{openproblem}

\begin{openproblem}[Four-dimensional gauge theory]
Construct four-dimensional Yang--Mills with mass gap as a K-S functor.  This
is at least as hard as the corresponding OS, Wightman, and Haag--Kastler
existence problems and additionally requires functorial gluing and
holomorphic admissible-metric dependence.
\end{openproblem}

\begin{openproblem}[Fermions and gauge fields in K-S form]
Develop K-S functorial constructions for fermionic and gauge theories in a
regularization that preserves the relevant reflection positivity, anomaly
bookkeeping, and sewing identities.  For gauge theory the BV formulation of
regularized path integrals and the lattice formulation are the currently
available frameworks with a precise gauge-consistency condition.
\end{openproblem}

\begin{openproblem}[Extended K-S structures]
Formulate and construct extended K-S theories assigning data to corners and
higher-codimension strata, analogous to extended TQFT but with non-topological
metric dependence and analytic sewing.  The target should record defects,
boundary conditions, and local operator categories.
\end{openproblem}
